\subsection{Additional Visual Results}

Рисунок~\ref{fig:app-exp01} показывает распределение разрывов между
структурной границей, оптимумом и исходной политикой. Нулевой разрыв означает
достижимость $B_4$, а не отсутствие очереди.

\begin{figure}[H]
  \centering
  \includegraphics[width=0.92\textwidth]{exp01_gaps.png}
  \caption{EXP-01: разрывы $T^*/B_4-1$ и разрывы исходной политики.}
  \label{fig:app-exp01}
\end{figure}

\begin{figure}[H]
  \centering
  \includegraphics[width=0.92\textwidth]{exp03_profile_effects.png}
  \caption{EXP-03: изменение $T^*$ при одинаковых агрегатах и разных фазовых
  профилях. Историческая метка sync здесь не означает синхронизацию
  wall-clock окон.}
  \label{fig:app-exp03}
\end{figure}

\begin{figure}[H]
  \centering
  \includegraphics[width=0.92\textwidth]{exp05_interaction_contrast.png}
  \caption{EXP-05: взаимодействие доступного параллелизма и гранулярности на
  канонической точной сетке.}
  \label{fig:app-exp05}
\end{figure}

\begin{figure}[H]
  \centering
  \includegraphics[width=0.90\textwidth]{candidate_task_lever_map.png}
  \caption{Иллюстративная локальная чувствительность варианта 4. В ячейках
  показано сокращение доказанного $T^*$ после контролируемого уменьшения фаз
  одной задачи; звёздочкой отмечены задачи номинального критического пути.
  Масштабы вмешательств синтетические и не являются эмпирическими оценками.}
  \label{fig:app-task-levers}
\end{figure}

\begin{figure}[H]
  \centering
  \includegraphics[width=0.92\textwidth]{exp06_ranking_stability.png}
  \caption{EXP-06: устойчивость ранжирования иллюстративных вариантов внутри
  замороженных семейств ошибок.}
  \label{fig:app-exp06}
\end{figure}

\begin{figure}[H]
  \centering
  \includegraphics[width=0.92\textwidth]{exp07_claim_effects.png}
  \caption{EXP-07: эффекты на синтетических DAG. C1 и C4 не поддержаны; C2 и
  C3 относятся преимущественно к фиксированной политике.}
  \label{fig:app-exp07}
\end{figure}

\begin{figure}[H]
  \centering
  \includegraphics[width=0.92\textwidth]{exp08_critical_path_effects.png}
  \caption{EXP-08: изменение заранее выбранного пути при $\gamma>C$ и
  переключения глобального критического пути.}
  \label{fig:app-exp08}
\end{figure}

\FloatBarrier

\subsection{Reproducibility}

Исходный код, замороженные входы и сохранённые результаты находятся в каталоге
\texttt{experiments}. Из корня проекта полный набор команд имеет вид:
\begin{verbatim}
cd article/experiments
uv sync --python 3.13
uv run pytest -q
uv run python -m experiments validate scenarios/article
uv run python -m experiments run --experiment EXP-00
uv run python -m experiments run --experiment EXP-01
uv run python -m experiments run --experiment EXP-02
uv run python -m experiments run --experiment EXP-03
uv run python -m experiments run --experiment EXP-04
uv run python -m experiments run --experiment EXP-05
uv run python -m experiments run --experiment EXP-06
uv run python -m experiments run --experiment EXP-07
uv run python -m experiments run --experiment EXP-08 --stage development
uv run python -m experiments run --experiment EXP-08 --stage confirmatory
uv run python -m experiments figures --set candidates
\end{verbatim}
EXP-00--EXP-07 независимы и могут повторяться по отдельности. Для EXP-08
порядок существенен: development создаёт \texttt{exp08\_freeze.json}, а
confirmatory-команда отказывается работать без совпадающего freeze и
замороженной матрицы. Shortcut \texttt{--stage all} выполняет те же две стадии
последовательно, но явные команды выше делают границу данных проверяемой.

\begin{center}
\begin{tabularx}{\textwidth}{@{}>{\raggedright\arraybackslash}p{0.39\textwidth}>{\raggedright\arraybackslash}X@{}}
\toprule
Путь & Содержимое \\
\midrule
\path{scenarios/article/} & полностью заданные иллюстративные варианты EXP-00 \\
\path{scenarios/canonical/} & замороженные матрицы EXP-01--EXP-08, исходный пилот EXP-01 и запись его глобального сокращения \\
\path{src/experiments/} & schema, validator, fixed policy, exact solver, генераторы серий и код рисунков \\
\path{tests/} & проверки схемы, генерации серий, инвариантов и регрессии чисел статьи \\
\path{results/exp??_manifest.json} & SHA-256 входов, реализации и выходов, версии зависимостей, settings и статусы \\
\path{results/exp??_report.md} & полный дизайн, таблицы результатов, ограничения и ссылки на raw CSV \\
\path{results/full_experiments_report.md} & единая сводка EXP-00--EXP-08 и карта артефактов \\
\path{figures/} & пары PNG/SVG и код их построения \\
\bottomrule
\end{tabularx}
\end{center}

Основные raw-таблицы хранят параметры сценария, $B_4$, срок фиксированной
политики, solver status, objective, best bound и gap; отдельные phase/event
таблицы позволяют повторно проверить ресурсную допустимость. Manifests
содержат commit Git и контрольные суммы перечисленных входов, реализации и
выходов, поэтому изменение любого перечисленного файла обнаруживается до
сопоставления результатов. Lockfile фиксирует
пакеты; сведения об ОС и hardware исходных прогонов в manifests отсутствуют,
что ограничивает воспроизводимость wall time, но не критерий
\texttt{OPTIMAL} и сохранённые целочисленные objectives.

\enlargethispage{2\baselineskip}
\paragraph{Data and code availability.}
Код лицензирован по MIT, а текст и исследовательские материалы --- по CC BY
4.0. Публичный репозиторий, версионированный релиз и DOI будут указаны после
архивирования версии препринта; состав дополнительного артефакта описан в
Разделе~\ref{sec:availability}.
